% ============================================================
%  DEFENSA DEL TRABAJO FIN DE GRADO (SECCIONES 7-9)
% ============================================================

% --- 7. JUSTIFICACIÓN Y ORIGINALIDAD DEL TFG
\section{Justificación y originalidad del TFG}

El dictamen que afirma que “El TFG presentado no está debidamente justificado y
carece de originalidad” puede rebatirse con tres argumentos complementarios:
metodológico, empírico y de aporte teórico. En primer lugar, desde la perspectiva
metodológica, el trabajo parte de una pregunta clara y relevante —la relación entre
los tiempos de espera en el Sistema Nacional de Salud y el gasto en seguros
sanitarios privados— y diseña un conjunto de especificaciones econométricas para
ponerla a prueba. La elección de un panel regional permite explotar la variación
temporal y espacial, controlar heterogeneza inobservable mediante efectos fijos
y estimar dinámicas mediante retardos y modelos umbral. Además, las decisiones
sobre deflactores, manejo de datos perdidos y robustez de errores estándar están
documentadas y justificadas en la memoria, lo que asegura la reproducibilidad.

En segundo lugar, desde el punto de vista empírico, el trabajo incorpora fuentes
diversas y actualizadas, combina series administrativas con encuestas y aplica
procedimientos de auditoría de datos que reducen sesgos de medición. No es un
ejercicio meramente descriptivo: se contrastan hipótesis específicas (elasticidad
ingreso, efecto umbral por tiempos extremos de espera, persistencia del gasto)
y se muestran resultados robustos a distintas especificaciones y muestras.

Finalmente, en términos de originalidad, el TFG aporta valor por su enfoque
combinado: integra la literatura sobre demanda de servicios privados con la de
efectos de oferta pública (listas de espera) y propone un umbral práctico que
captura no linealidades relevantes para la política pública. El uso de modelos
umbral aplicados a un panel de CCAA y la traducción de efectos logarítmicos a
magnitudes económicas interpretables constituyen contribuciones novedosas a
nivel de trabajo fin de grado. En suma, la justificación teórica y empírica es
coherente, reproducible y con aportes concretos que contradicen el dictamen
critico planteado.

% --- 8. CONTRIBUCIÓN DEL TFG
\section{Contribución del TFG}

Contrarrestar la aseveración “El TFG presentado no aporta ni supone contribución
alguna a su formación” exige distinguir entre contribución académica y formativa.
Académicamente, el TFG articula hipótesis contrastables y produce evidencias que
enriquecen la discusión sobre incentivos al gasto privado en sanidad. La estimación
de elasticidades y la identificación de un umbral en tiempos de espera ofrecen
información práctica para gestores y decisores que no aparece de forma tan
explícita en trabajos previos a nivel autonómico; además, el estudio realiza
pruebas de robustez y sensibilidad que fortalecen la validez de las conclusiones.

En lo formativo, la contribución es clara y multidimensional: el autor ha
aplicado técnicas econométricas avanzadas (modelos de panel dinámicos, tests de
Hausman, errores agrupados) y ha manejado bases de datos reales, limpiado series
y elaborado tablas y visualizaciones para sacar implicaciones políticas. Este
proceso desarrolló habilidades prácticas —programación en Python/R, manejo de
datos panel, interpretación estadística y comunicación de resultados— que son
centrales para la formación profesional en economía de la salud.

Además, el documento propone recomendaciones de política concretas (por
ejemplo, acciones focalizadas para reducir esperas extremas y su evaluación
ex post) y líneas futuras de investigación (instrumentos para endogeneidad,
extensión a microdatos), mostrando que el trabajo no solo resume conocimientos
sino que abre vías de continuación, lo que es un indicador claro de aportación
intelectual. En conjunto, estas dimensiones demuestran que el TFG aporta valor
tanto al cuerpo académico como a la formación práctica del autor.

% --- 9. VALORACIÓN CRÍTICA DEL TFG
\section{Valoración crítica del TFG}

Un análisis honesto de fortalezas y debilidades es esencial para la defensa.
Entre las fortalezas, destaco la claridad de la pregunta de investigación y la
consistencia del diseño empírico: la combinación de modelos de panel, análisis
dinámico y especificaciones de umbral permite capturar múltiples canales de
efecto. La documentación de las decisiones sobre tratamiento de datos y la
serie de pruebas de robustez incrementan la credibilidad de los resultados.

En cuanto a debilidades, la principal limitación es la naturaleza agregada de
las unidades de análisis (comunidades autónomas), que impide observar heterogeneidades
individuales (por ejemplo, diferencias socioeconómicas dentro de CCAA) y limita
la identificación causal en presencia de shocks regionales no observados. Otra
debilidad es la potencial endogeneidad entre gasto privado y factores institucionales
no observados que podrían influir simultáneamente en las listas de espera; aunque
se proponen instrumentos y pruebas de sensibilidad, estos enfoques requieren
datos adicionales para una validación concluyente.

Finalmente, en términos de alcance, el TFG prioriza una primera aproximación
empírica sólida frente a la exploración de mecanismos microeconómicos detallados;
por tanto, una extensión natural sería trabajar con microdatos o con experimentos
naturales para fortalecer la inferencia causal. Estas limitaciones son, sin
embargo, gestionables y abren oportunidades claras de trabajo futuro.

% Fin del capítulo de defensa (secciones 7-9)
